\documentclass{article}
\usepackage{fontspec}
\usepackage{type1cm}
\usepackage{geometry}
\usepackage[bold-style=ISO]{unicode-math}
\usepackage[heading=true]{ctex}%添加heading=true，使用中文版式
\usepackage{graphicx}
\geometry{a4paper,left=1cm,right=1cm,top=1cm,bottom=1cm}
\begin{document}
\begin{flushleft}
	\fontsize{24pt}{30pt}\selectfont
	~\\ \textbf{导数} \\~\\
	可导必连续,不连续必不可导\\
	~\\
	求某一点$x_0$的导数:用导数的定义$\lim\limits_{x\to x_0} \frac{f(x)-f(x_0
		)}{x-x_0}$\\
	~\\
	导数的几何意义:切线的斜率\\
	切线方程:$y-y_0=k(x-x_0)$\\
	~\\
	设$f(x)$是$(-a,a)$上的偶(奇)函数且可导,则$f(x)^{'}$是$(-a,a)$上的奇(偶)函数\\
	设$f(x)$以$T$为周期且可导,则$f(x)^{'}$也以$T$为周期\\
	~\\
	可导$\pm$不可导$=$不可导\\
	~\\
	$y=f(x)$的反函数为$x=\phi(y)$,则$f^{'}(x)=\frac{1}{\phi^{'}(y)}$\\
	~\\
	复合函数求导:从外到里层层求导\\
	~\\
	$f(x)$可导,$g(x)$连续但不可导,若$f(x)\times g(x)$在$x_0$处可导,则$f(x_0)=0$\\
	~\\ \textbf{基本导数公式}\\~\\
	$(C)^{'}=0$\\
	$(x^u)^{'}=ux^{u-1}$\\
	$(\sin x)^{'}=\cos x$\\
	$(\cos x)^{'}=-\sin x$\\
	$(\tan x)^{'}=\sec^2 x$\\
	$(\cot x)^{'}=-\csc^2 x$\\
	$(\sec x)^{'}=\sec x\tan x$\\
	$(\csc x)^{'}=-\csc x\cot x$\\
	$(a^x)^{'}=a^xlna$\\
	$(e^x)^{'}=e^x$\\
	$(log_ax)^{'}=\frac{1}{xlna}$\\
	$(lnx)^{'}=\frac{1}{x}$\\
	$(\arcsin x)^{'}=\frac{1}{\sqrt{1-x^2}}$\\
	$(\arccos x)^{'}=-\frac{1}{\sqrt{1-x^2}}$\\
	$(\arctan x)^{'}=\frac{1}{1+x^2}$\\
	$($arccot$ x)^{'}=-\frac{1}{1+x^2}$\\
	$(ln|x|)^{'}=\frac{1}{x}$\\
	~\\ \textbf{导数的四则运算}\\~\\
	设$u,v$均可导,则:\\
	$(u\pm v)^{'}=u^{'}+v^{'}$\\
	$(uv)^{'}=u^{'}v+uv^{'}$\\
	$(\frac{u}{v})^{'}=\frac{u^{'}v-uv^{'}}{v^2}$\\
	~\\ \textbf{求n阶导数}\\~\\
	$(e^{ax+b})^{(n)}=a^ne^{ax+b}$\\
	$[\sin(ax+b)]^{(n)}=a^n\sin(ax+b+\frac{n\pi}{2})$\\
	$[\cos(ax+b)]^{(n)}=a^n\cos(ax+b+\frac{n\pi}{2})$\\
	$[ln(ax+b)]^{(n)}=(-1)^{n-1}a^n\frac{(n-1)!}{(ax+b)^n}$\\
	$(\frac{1}{ax+b})^{(n)}=(-1)^{n}a^n\frac{(n)!}{(ax+b)^{n+1}}$\\
	~\\
	若$f(x)$可分解成$f_1(x)$和$f_2(x)$,且知道$f_1^{(n)}(x)$和$f_2^{(n)}(x)$,则$f^{(n)}(x)=f_1^{(n)}(x)+f_2^{(n)}(x)$\\
	~\\
	莱布尼兹公式:\\
	$(uv)^{(n)} = \sum\limits_{k=0}^{n} C_n^k u^{(n-k)} v^{(k)} = \sum\limits_{k=0}^{n} C_n^k u^{(k)} v^{(n-k)}$\\
	$C_n^m=\frac{n!}{m!(n-m)!}$\\
	$C_n^0=1$\\
	~\\ \textbf{隐函数求导}\\~\\
	$y=y(x)$由方程$F(x,y)=0$确定,\\
	在方程$F(x,y)=0$两端直接对$x$求导\\
	~\\ \textbf{参数方程求导}\\~\\
	$y=y(x)$由参数方程$\left\{
	\begin{array}{rcl}
	x=\phi(t)\\
	y=\Phi(t)
	\end{array} \right.$确定,且$\phi(t)$和$\Phi(t)$均可导,则$y^{'}=\frac{y^{'}_t}{x^{'}_t}=\frac{\Phi^{'}(t)}{\phi^{'}(t)}$\\
	~\\ \textbf{微分}\\~\\
	微分的定义:设函数$y=f(x)$在$x_0$的邻域内有定义,如果函数的增量$\Delta y = f(x_0 + \Delta x) − f(x_0)$可表示为$ \Delta y = A\Delta x + o(\Delta x)$(其中$A$是不依赖于$\Delta x$的常数,那么称函数$f(x)$在点$x_0$是可微的,且$A\Delta x$称作函数在点$x_0$相应于自变量增量$\Delta x$的微分，记作$dy$，即$dy = A\Delta x$\\
	若$x$是自变量,则$\Delta x=dx$\\
	~\\
	可微必可导,可导必可微,且$dy=f^{'}(x)dx$\\
	~\\
	\includegraphics[scale=1.0]{1.jpg}
	
\end{flushleft}
\end{document}